\chapter{Tapentadol}
\index{Tapentadol}

\medskip
\noindent\textit{Educational reference; always seek professional advice for individual care.}

\section*{Definition and Overview}
Tapentadol is a prescription medicine for moderate to severe pain. It belongs to a class of medicines called opioid analgesics, but differs from most other opioids because it works in two distinct ways at once. Tapentadol is available in many countries under the brand name Palexia, in immediate-release and prolonged-release forms, and also as a generic medicine. It is a Schedule 8 (controlled drug) medicine in most local states and territories, meaning it is tightly regulated and prescribed only by medical practitioners because of its potential for dependence and misuse. Tapentadol is used for pain severe enough to require a strong pain reliever, such as pain after surgery, pain from injuries, and certain types of chronic pain including musculoskeletal and neuropathic pain. It is sometimes chosen because it may cause less nausea and constipation than some other opioids, although individual responses vary greatly. Because tapentadol is a powerful medicine with a real risk of serious harm if used incorrectly, it should only ever be taken exactly as prescribed. This chapter explains what tapentadol is, how it works, the important adverse effects and interactions to watch for, and the practical steps that keep its use safe.

\section*{Epidemiology}
Opioid medicines are among the most commonly prescribed strong analgesics worldwide, and tapentadol is one of the newer members of this group. It was introduced onto the local market in the early 2010s and has since become widely used, particularly in chronic non-cancer pain and in older adults, in whom its side-effect profile is sometimes preferred. It is now one of the more commonly dispensed strong opioids in many countries, though it is far less commonly prescribed than codeine and oxycodone. Use is more common among middle-aged and older people than among young people, reflecting the age pattern of the pain conditions it treats, and it is used about equally in men and women. As with all opioids, there is ongoing concern about long-term use, dose escalation, and dependence, and local data have shown that a small but important number of people who start tapentadol for chronic pain continue taking it for long periods. The emphasis in local prescribing has shifted firmly towards short courses, the lowest effective dose, and careful monitoring of anyone who needs it long term.

\section*{Aetiology and Pathophysiology}
Tapentadol works on the central nervous system through two complementary mechanisms that together reduce the perception of pain.

\begin{itemize}
    \item \textbf{Mu-opioid receptor agonism:} tapentadol activates the same opioid receptors in the brain and spinal cord that are triggered by morphine and other opioids, reducing the transmission of pain signals and the emotional response to pain
    \item \textbf{Noradrenaline reuptake inhibition:} tapentadol also slows the removal of noradrenaline from the spaces between nerve cells, increasing the activity of descending pain-inhibiting pathways that travel from the brain down the spinal cord, an effect shared with some antidepressants used for nerve pain
    \item \textbf{Rapid absorption and elimination:} the medicine is well absorbed after swallowing, immediate-release tablets begin working within about an hour, and the active drug is cleared by the kidneys, with dosing adjusted in significant kidney or liver disease
    \item \textbf{Minimal enzyme interactions:} unlike many opioids, tapentadol is metabolised largely by conjugation rather than the cytochrome P450 system, giving it fewer interactions with other medicines
\end{itemize}

The dual mechanism is thought to explain why tapentadol can relieve both pain from tissue damage and pain from nerve damage, and why some people find it causes less nausea, drowsiness, and constipation than morphine or oxycodone. The downside is that the medicine can still cause the characteristic harms of the opioid group, including dependence, tolerance, withdrawal, and breathing suppression, and the noradrenaline effect adds a small risk of blood pressure changes and serotonin-type side effects.

\section*{Clinical Features}
Tapentadol produces both the intended relief of pain and a predictable range of unwanted effects. Features people most commonly notice include:

\begin{itemize}
    \item \textbf{Pain relief:} the intended effect, which usually begins within an hour for immediate-release forms and builds gradually for prolonged-release forms taken once or twice a day
    \item \textbf{Drowsiness and sedation:} very common, especially in the first days of treatment, and a major reason for avoiding driving and operating machinery until the response is known
    \item \textbf{Dizziness and light-headedness:} common, especially when standing up quickly, and more likely in older people
    \item \textbf{Nausea and vomiting:} common early in treatment, often settling over the first week, and usually controlled with anti-nausea medicines if needed
    \item \textbf{Constipation:} frequent with all opioids, including tapentadol, and best prevented with fluids, fibre, and activity
    \item \textbf{Low blood pressure:} a drop in blood pressure on standing, which can cause fainting and falls in older people
    \item \textbf{Slowed breathing:} in higher doses, or when combined with alcohol or other sedatives, tapentadol can slow breathing, the most dangerous of its effects
    \item \textbf{Withdrawal symptoms:} restlessness, sweating, muscle aches, runny nose, cramps, and agitation can occur if the medicine is stopped suddenly after regular use
\end{itemize}

\section*{Differential Diagnosis}
Because tapentadol's adverse effects overlap with many common conditions, a doctor assessing someone taking the medicine must distinguish between drug effects and other causes.

\begin{itemize}
    \item \textbf{Constipation from other causes:} a low-fibre diet, dehydration, or other medicines such as antidepressants and iron can all cause constipation wrongly blamed on tapentadol
    \item \textbf{Nausea from a gastrointestinal illness:} gastroenteritis or reflux can cause nausea and vomiting in someone taking tapentadol, and a careful history usually separates the causes
    \item \textbf{Dizziness from inner-ear or neurological disease:} vertigo, labyrinthitis, or a minor stroke can mimic the light-headedness of the medicine, especially in older people
    \item \textbf{Depression and fatigue:} the sedation of opioids can be confused with a mood disorder, particularly in people with low mood and poor sleep
    \item \textbf{Serotonin syndrome versus opioid effects:} confusion, agitation, sweating, tremor, and a fast heart rate can arise from the interaction of tapentadol with antidepressants, and this needs urgent assessment
    \item \textbf{Pain from worsening disease:} if the original condition is progressing, increasing pain may be mistaken for the medicine stopping working rather than a need to reassess the diagnosis
\end{itemize}

\section*{When to Seek Medical Care}
Anyone prescribed tapentadol should have a clear plan for when to seek review. Contact your doctor or pharmacist if you have pain that is not controlled, if you feel unusually drowsy, if you become constipated for more than a few days, or if you are pregnant, planning pregnancy, or breastfeeding and have not discussed this. Do not stop the medicine suddenly, and never change your dose without advice.

\textbf{Contact your local emergency services for:}
\begin{itemize}
    \item Severe drowsiness, difficulty staying awake, slow or shallow breathing, or noisy breathing
    \item Confusion, agitation, hallucinations, or a very fast or irregular heart rate, which can suggest serotonin syndrome
    \item Fainting, collapse, or a seizure
    \item Swelling of the face, lips, or tongue, or difficulty breathing
    \item Suspected overdose, or a child or another person accidentally taking the medicine
\end{itemize}

\section*{Diagnosis and Investigations}
Tapentadol is a prescription medicine, so the most important diagnostic step is a careful medical review rather than laboratory testing.

\begin{itemize}
    \item \textbf{Medication review:} a doctor or pharmacist reviews the full medicine list, including over-the-counter and complementary medicines, to check for interactions
    \item \textbf{Kidney and liver tests:} blood tests may be ordered if there is known kidney or liver disease, because these organs clear the medicine from the body
    \item \textbf{Blood pressure measurement:} checking lying and standing blood pressure helps detect the drop that can cause falls
    \item \textbf{Pain assessment:} the doctor reviews the type, site, and severity of pain and how well the medicine is controlling it
    \item \textbf{Screening for misuse risk:} doctors may use questionnaires and check the prescription database (such as SafeScript and similar systems) to detect multiple prescribers or unusual use
    \item \textbf{Driving and pregnancy review:} drowsiness affects driving safety, and in pregnancy or breastfeeding the benefits and risks are weighed carefully
\end{itemize}

\section*{Management}
The safe use of tapentadol depends on the right dose, the right duration, and careful review.

\begin{itemize}
    \item \textbf{Lowest effective dose:} the starting dose is low and adjusted according to response, with the goal of using the smallest dose that controls pain
    \item \textbf{Shortest effective course:} for acute pain, the medicine is used for days to weeks and then stopped, usually with a tapering plan
    \item \textbf{Regular timing:} immediate-release tablets are taken when needed up to a prescribed maximum, while prolonged-release tablets are taken at the same time each day and should not be crushed or chewed
    \item \textbf{Non-drug pain measures:} physiotherapy, heat or cold, gentle activity, pacing, and psychological approaches are used alongside the medicine, especially for chronic pain
    \item \textbf{Other pain medicines:} paracetamol and some anti-inflammatories may be combined with tapentadol under medical advice to improve relief without raising the opioid dose
    \item \textbf{Preventing constipation and nausea:} fluids, fibre, regular activity, a laxative if needed, taking doses with food, and anti-nausea medicines in the first week
    \item \textbf{Supervised tapering:} when the medicine is no longer needed, the dose is reduced gradually to avoid withdrawal symptoms
    \item \textbf{Naloxone awareness:} in people at higher risk of overdose, doctors may prescribe naloxone and teach family members how to use it as an emergency antidote
    \item \textbf{Referral:} people with chronic pain are often referred to a pain clinic, psychologist, or rehabilitation service for a broader plan
\end{itemize}

\section*{Complications}
The complications of tapentadol are those of any strong opioid, and most are preventable with careful prescribing.

\begin{itemize}
    \item \textbf{Tolerance:} the same dose becomes less effective over time, which can tempt people to take more than prescribed
    \item \textbf{Dependence:} physical and psychological dependence can develop with regular use, making it hard to stop without support
    \item \textbf{Withdrawal:} stopping suddenly causes agitation, sweating, muscle aches, cramps, and insomnia, which is why doses are tapered
    \item \textbf{Respiratory depression:} slow or shallow breathing, most dangerous in overdose or when combined with alcohol or sedatives
    \item \textbf{Serotonin syndrome:} a potentially life-threatening reaction when tapentadol is combined with serotonergic medicines such as SSRI antidepressants
    \item \textbf{Falls and fractures:} drowsiness, dizziness, and low blood pressure increase the risk of falls, especially in older people
    \item \textbf{Severe constipation:} can occasionally progress to a blocked bowel
    \item \textbf{Overdose:} taking too much, or combining the medicine with alcohol or other sedatives, can cause unconsciousness and fatal breathing failure
\end{itemize}

\section*{Prognosis and Outcomes}
For acute pain, such as pain after an operation or a fracture, tapentadol is highly effective for the first days to weeks, and most people can stop it within a short time without difficulty. The prognosis is excellent when the course is short and the dose modest. For chronic pain the picture is different: no opioid cures the underlying condition, and long-term use brings modest benefit to some people and real harm to others. Studies of tapentadol in chronic musculoskeletal pain show that a substantial minority gain useful relief, but tolerance, adverse effects, and dependence can limit its value over time, which is why local guidelines recommend regular review, the lowest effective dose, and a plan for reducing or stopping the medicine if the pain does not improve. Most people who experience adverse effects find they settle with dose adjustment or a change of medicine. With careful medical supervision, tapentadol can be a useful part of a pain-management plan, but it works best combined with non-drug treatments and an honest discussion of goals and expectations.

\section*{Prevention and Self-Care}
\begin{itemize}
    \item Take tapentadol exactly as prescribed, at the stated times, and never exceed the dose
    \item Store the medicine in its original pack in a secure place, out of sight and reach of children and visitors
    \item Keep one current, up-to-date list of all your medicines and show it to your doctor and pharmacist at every visit
    \item Do not drink alcohol or take sleeping tablets, anxiety medicines, or other sedatives without checking with your doctor or pharmacist
    \item Ask your pharmacist about every new medicine, including herbal and complementary products, before combining it with tapentadol
    \item Return unused or expired tablets to a pharmacy for safe disposal, and never share the medicine with anyone else
    \item Report drowsiness, constipation, or falls to your doctor promptly rather than tolerating them
    \item If you are stopping the medicine, follow your doctor's tapering plan and never stop abruptly
    \item Talk to your doctor if the medicine no longer controls pain, and avoid silently increasing your own dose
\end{itemize}

\section*{Sources and Further Reading}
The following organisations provide reliable information about tapentadol and the safe use of opioid pain medicines.

\begin{itemize}
    \item healthdirect Australia. \textit{Tapentadol: consumer medicine information.} https://www.healthdirect.gov.au/
    \item Therapeutic Goods Administration. \textit{Medicines and scheduling information for tapentadol.} https://www.tga.gov.au/
    \item NPS MedicineWise. \textit{Opioid medicines: reducing harms, using wisely.} https://www.nps.org.au/
    \item NHS. \textit{Opioid pain medicines: uses, cautions, and adverse effects.} https://www.nhs.uk/
    \item Mayo Clinic. \textit{Pain management: using opioid medicines safely.} https://www.mayoclinic.org/
    \item Australian Department of Health and Aged Care. \textit{Chronic pain and opioid use resources.} https://www.health.gov.au/
\end{itemize}

\clearpage
